\documentclass[11pt]{article}
\usepackage[margin=1in]{geometry}
\usepackage{hyperref}
\usepackage{enumitem}
\usepackage{amsmath,amssymb,amsthm,mathtools}
\usepackage[T1]{fontenc}
\usepackage[utf8]{inputenc}
\title{\textbf{Theory-mode report: Fluctuations for the Toda lattice}}
\author{Generated from Scholar Inbox digest}
\date{Digest date: 2026-04-17}
\begin{document}
\maketitle
\section*{Paper information}
\begin{itemize}[leftmargin=1.5em]
\item \textbf{Title:} Fluctuations for the Toda lattice
\item \textbf{Authors:} Amol Aggarwal, Matthew Nicoletti
\item \textbf{arXiv:} \href{https://arxiv.org/abs/2604.14346}{2604.14346}
\item \textbf{Scholar digest score:} 77.0
\end{itemize}
\section*{Executive summary}
This report summarizes the highest-scoring paper in the Scholar Inbox digest dated \textbf{2026-04-17}. The abstract states the core claim as follows: \emph{In this paper we consider the Toda lattice \$(\textbackslash\{\}mathbf\{p\}(t);\textbackslash\{\}mathbf\{q\}(t))\$ at thermal equilibrium, meaning that its variables \$(p\_j)\$ and \$(e\textasciicircum{}\{q\_j - q\_\{j+1\}\})\$ are independent Gaussian and Gamma random variables, respectively. We show under diffusive scaling that the space-time fluctuations for the model's currents converge to an explicit Gaussian limit. As consequences, we deduce, (i) the scaling limit for the trajectory of a single particle \$q\_0\$ is a Brownian motion; (ii) space-time two-point correlation functions for the model decay inversely with time, with explicit scaling distributions predicted by Spohn (Spohn, J. Phys. A 53 (2020), 265004). Our starting point is the notion that the Toda lattice can be thought of as a dense collection of many ``quasi-particles'' that interact through scattering. The core of our work is to establish that the full joint scaling limit of the fluctua}.

\section*{Problem setup and intuition}
The introduction reinforces this lens: the paper appears motivated by a gap between practical behavior and theoretical understanding.

\section*{Proof architecture}
The technical core appears to build the surrogate quantity and the chain of lemmas needed to propagate local control into the final theorem: the paper follows a standard theory-paper structure with structural reduction plus quantitative bounds.

\section*{Takeaway}
The key reusable lesson is to define the right object, find the structural reduction, and quantify the error terms clearly.
\end{document}
